\chapter{Results and
Evaluation}
\label{chapter-10-results-and-evaluation}

\section{Evaluation Methodology}\label{evaluation-methodology}

The efficacy of the PhishGuard machine learning architecture was
rigorously evaluated using a holdout methodology to prevent data leakage
and ensure generalized performance on unseen data. The dataset,
comprising diverse phishing and legitimate Uniform Resource Locators
(URLs), was systematically split into training and testing partitions.

The primary training corpus consisted of 23,374 samples
(\texttt{trainSamples:\ 23,374}). To accurately assess the model's
predictive capability in a simulated real-world environment, a strictly
held-out test set comprising 5,009 samples
(\texttt{testSamples:\ 5,009}) was established. This test set maintained
an near-perfect class balance, containing 2,505 legitimate instances and
2,504 phishing instances, ensuring that evaluation metrics were not
artificially skewed by class imbalance.

\section{Hyperparameter
Optimization}\label{hyperparameter-optimization}

Prior to final evaluation, the XGBoost classifier underwent extensive
hyperparameter tuning utilizing the Optuna optimization framework. The
optimization process executed 50 distinct trials (\texttt{nTrials:\ 50})
employing 5-fold cross-validation (\texttt{nCvFolds:\ 5}) to iteratively
search the hyperparameter space for the configuration that maximized the
Area Under the Receiver Operating Characteristic Curve (ROC AUC).

The optimal configuration was identified at Trial 43
(\texttt{bestTrialNumber:\ 43}), yielding a cross-validated AUC of
0.9943. The resulting best hyperparameters significantly constrained the
model's complexity to prevent overfitting while maintaining high
predictive capacity: - \textbf{Maximum Depth (\texttt{max\_depth}):} 7 -
\textbf{Number of Estimators (\texttt{n\_estimators}):} 700 -
\textbf{Learning Rate (\texttt{learning\_rate}):} \textasciitilde0.177 -
\textbf{Subsample Ratio (\texttt{subsample}):} \textasciitilde0.945 -
\textbf{Column Subsample by Tree (\texttt{colsample\_bytree}):}
\textasciitilde0.873 - \textbf{Gamma (\texttt{gamma}):}
\textasciitilde0.198

\section{Empirical Performance
Metrics}\label{empirical-performance-metrics}

The fully optimized XGBoost model, utilizing the complete 21-dimensional
feature vector (17 URL structural features and 4 OSINT-derived
features), demonstrated exceptional performance on the held-out test set
of 5,009 samples.

The model achieved an aggregate \textbf{Accuracy of 96.45\%}
(\texttt{0.96446}), indicating robust overall correctness. However, in
the context of threat detection, precision and recall offer more nuanced
insights into operational viability.

The system achieved a \textbf{Precision of 97.86\%} (\texttt{0.97860}),
demonstrating a low false-positive rate---a critical metric for
minimizing alert fatigue among end-users. Correspondingly, the
\textbf{Recall reached 94.97\%} (\texttt{0.94968}), signifying the
model's strong capability to successfully identify true phishing threats
within the dataset. The harmonic mean of these two metrics resulted in
an \textbf{F1-Score of 96.39\%} (\texttt{0.96392}).

Furthermore, the model exhibited outstanding discriminative ability
across various classification thresholds, achieving a \textbf{ROC AUC of
99.41\%} (\texttt{0.99408}) and a Precision-Recall Area Under Curve
(\textbf{PR AUC}) of \textbf{99.48\%} (\texttt{0.99479}).

\subsection{Error Analysis and Confusion
Matrix}\label{error-analysis-and-confusion-matrix}

A granular analysis of the classification errors on the test set reveals
the model's operational tendencies. Out of 5,009 total predictions: -
\textbf{True Positives (Phishing correctly identified):} 2,378 -
\textbf{True Negatives (Legitimate correctly allowed):} 2,453 -
\textbf{False Positives (Legitimate incorrectly flagged):} 52 -
\textbf{False Negatives (Phishing incorrectly allowed):} 126

The disparity between False Positives (52) and False Negatives (126)
underscores the model's high-precision orientation. While the system is
highly reliable when it issues a warning, a small subset of
sophisticated phishing URLs successfully evaded the structural and OSINT
feature checks.

\section{Feature Explainability and SHAP
Analysis}\label{feature-explainability-and-shap-analysis}

To demystify the XGBoost model's decision-making process and move beyond
``black-box'' predictions, SHapley Additive exPlanations (SHAP) were
employed. The SHAP \texttt{TreeExplainer} was utilized to compute the
marginal contribution of each of the 21 features across the test set.

The SHAP analysis (\texttt{shapAnalysis.py}) revealed a distinct
hierarchy in feature importance. The presence of HTTPS
(\texttt{isHttps}) emerged as the most influential single feature
(accounting for roughly 33.4\% of the global importance), followed
closely by the validity of the domain's DNS configuration
(\texttt{hasValidDns}, 12.5\%). Structural anomalies, such as
\texttt{specialCharCount} (8.5\%) and \texttt{pathDepth} (7.4\%), also
demonstrated significant predictive power.

\subsection{OSINT Ablation Study}\label{osint-ablation-study}

A critical objective of the PhishGuard architecture was to evaluate the
supplemental value of Open-Source Intelligence (OSINT) data when
combined with traditional URL lexical analysis. To quantify this, an
ablation study was conducted.

The global SHAP contribution analysis (\texttt{ablation\_report.json})
calculated the mean absolute SHAP values for the two distinct feature
subsets. The analysis determined that the 17 URL structural features
account for \textbf{85.64\%} of the model's total predictive influence.
The 4 dynamic OSINT features (\texttt{usesCdn}, \texttt{dnsRecordCount},
\texttt{hasValidDns}, \texttt{hasValidMx}) contributed the remaining
\textbf{14.36\%}.

Interestingly, when the model was retrained exclusively on the 17 URL
features (excluding all OSINT data), the baseline performance metrics
exhibited a nominal, fractional increase (Accuracy shifted by
\textasciitilde+0.53\%). This phenomenon suggests that within this
specific, balanced dataset, the structural URL anomalies were
sufficiently profound to classify the targets independently. The
inclusion of OSINT features, while contributing 14.36\% to the model's
SHAP explanations and providing critical human-readable context for the
heuristic Scorer module, acted slightly as a regularizer in the pure ML
context, mitigating over-reliance on purely lexical characteristics.

